\documentclass[Total.tex]{subfiles}

\begin{document}
	
	\chapter{Bundles}
	
	\section{Bundles}
	What is bundle?
	
	\subsection{Definitions}
	
	\begin{Def}
		\begin{itemize}
			\item In \underline{most} general case, a \textbf{bundle} over an object $B$ in a category $\mathcal{C}$ is simply an object $E$ of $\mathcal{C}$ equipped wiht a morphism in $\mathcal{C}$ from $E$ to $B$
			
			\begin{gather*}
				p:E\rightarrow B
			\end{gather*}
			
			And let $(x:U\rightarrow B)$ be a \underline{generalized element} of $B$. (\ref{Category-Theory-Generalized-Element}), the \textbf{fibre} $E_x$ of the bundle over $x$ is the pullback is the pullback $x^*E$
			\begin{center}
				\begin{tikzcd}
					x^*E=U\times_B E\arrow[d]\arrow[r]&E\arrow[d,"p"]\\
					U\arrow[r,"x"]&B
				\end{tikzcd}
			\end{center}
			Here, $U=1$.
			
			Such bundles over $B$ is a category. Called \textbf{slide category} $\mathcal{C}_{/B}$ of $\mathcal{C}$ over $B$.
			
			\item The morphism $E_1\rightarrow E_2$ is a diagram such that the following diagram commutes
			\begin{center}
				\begin{tikzcd}
					E_1\arrow[rr]\arrow[rd]&&E_2\arrow[ld]\\
					&B
				\end{tikzcd}
			\end{center}
		\end{itemize}
	\end{Def}
	
	Some of Bundles are listed as following:
	
	\begin{itemize}
		\item Fibre Bundle;
		\item Micro Bundle;
		\item Principle Bundle;
		\item Vector Bundle;
		\item Bundle with Connection;
	\end{itemize}
	
	\subsubsection{Cross Sections}
	\begin{Def}
		A \textbf{cross section} of a bundle $(E,p,B)$ is a 
	\end{Def}
	
	\subsubsection{Product Bundle}
	
	\subsubsection{Stiefel Variety}
	
	\begin{Def}
		内容...
	\end{Def}
	
	\subsubsection{Grasmannian Variety}
	\begin{Def}
		内容...
	\end{Def}
	\subsection{Morphsims of Bundles}
	\subsection{Product and Fibre Products}
	\subsection{Restriction and Induced Bundles}
	\subsection{Local Properties}
	
	\section{Connections}
	
	\chapter{Fibre Bundles}
		\section{Fibre Bundles}
			\begin{Def}
				Let $\eta:E\rightarrow B$ be a bundle. Then, a \textbf{fiber bundle} is $(\eta,E,B,F)$ with standard fiber $F$, such that for every general element $1\rightarrow B$, the pullback of $E$ along $x$ is isomorphic to $F$.
				
				i.e. $E_x\cong F$.
			\end{Def}
	\chapter{Associated Bundles}
	\chapter{Projective Modules and Vector Bundles}
	\section{Free module, $GL_n$ and Stably Free Modules}
	
	\subsection{The Module over Rings}
	
	Review:
	
	\begin{itemize}
		\item Vector Space(over division ring);
		
		\item Base;
		
		\item IBP property;
		
		\begin{lemma}
			If there exists a ring map $f:R\rightarrow F$, where $F$ is a field or division ring, then $R$ has IBP property.
		\end{lemma}
		
		\begin{proof}
			Consider the following map
			
			\begin{gather*}
				M\rightarrow M\otimes_R F=V\qquad m\mapsto m\otimes_R 1
			\end{gather*}
			
			Then, $\sum \lambda_i m_i=0\rightarrow \lambda_i=0\Leftrightarrow \sum \lambda_i (m_i\otimes_R 1)=0\rightarrow \lambda_i=0$.
			
			Especially, if commutative, we could take $F=R/m$.
		\end{proof}
		
		Thus, $dim(V)=dim(M)$ is well-defined.
		
		\begin{Def}
			In general, for the ring $R$, take the notation of 
			
			\begin{gather*}
				GL_n(R)
			\end{gather*}
			
			for the group of invertible matrices. Especially, $n=1$, $GL_1(R)=R^\times$.
		\end{Def}
		
		\begin{lemma}
			Any finite-dimensional algebra over a field/division ring must satisfies $IBP$, because the rank of free $R$-module $M$ is an invariant:
			
			\begin{gather*}
				rank(M)=dim_F(M)/dim_F(R)
			\end{gather*}
		\end{lemma}
		
		\begin{proof}
			内容...
		\end{proof}
		
	\end{itemize}
	
	\subsection{Stably Free}
	
	\begin{Def}
		An $R$-module $P$ is called \textbf{stably free} of rank $n-m$ if
		\begin{gather*}
			P\oplus R^m\cong R^n
		\end{gather*}
		for some $m,n$. 
		
		Conversely:
	\end{Def}
	
	\begin{example}
		内容...
	\end{example}
	
	
	
	\subsection{Stably Free Property}
	
	\begin{Def}
		内容...
	\end{Def}
	
	\subsection{Some of Noncommutative Ring Theory}
	
	\begin{proposition}
		(P.M.Cohn) Consider the following conditions on a ring $R$:
		
		\begin{itemize}
			\item $R$ satisfies IBP;
			\item For all $m,n$, if $R^m\cong R^n\oplus P$, then $m\geq n$;
			\item For all $n$, if $R^n\cong R^n\oplus P$, then $P=0$.
		\end{itemize}
	\end{proposition}
	
	\begin{proof}
		内容...
	\end{proof}
	
	\section{Projective Modules}
	
	Given in noncommutative ring theory/Homological Algebra
	
	\section{Picard Group of Commutative Rings}
	
	Given in noncommutative ring theory.
	
	\section{Topological Vector Bundle, and Chern Classes}
	
	We will work on field $F=\RR,\CC$ or $\mathbb{H}$.
	
	\subsection{Families of Vector Spaces}
	
	\begin{Def}
		\begin{itemize}
			\item Let $X$ be a topological space. A family of vector spaces is a topological space $E$, together with map
			\begin{gather*}
				\eta:E\rightarrow X
			\end{gather*}
			with a finite dimensional vector space structure on each $E_x=\eta^{-1}(\{x\})$. 
			
			The structure of $E_x$ is compatible with the topological structure of $E$. i.e.
			\begin{gather*}
				F\times E\rightarrow E\qquad E\times_X E\rightarrow E
			\end{gather*}
			are continuous.
			
			\item  A \textbf{homomorphism} from $\eta:E\rightarrow X$ to $\varphi:F\rightarrow X$ is:
			
			\begin{itemize}
				\item A continuous map $f:E\rightarrow F,\eta=\varphi f$;
				\item $f_x:E_x\rightarrow F_x$ is a linear map of vector spaces;
			\end{itemize}
		\end{itemize}
	\end{Def}
	
	
	
	If $Y\subseteq X$, then define
	
	\begin{gather*}
		E|Y=(E|Y,Y,\eta|\eta^{-1}(Y))
	\end{gather*}
	With restriction map $\eta|\eta^{-1}(Y):\eta^{-1}(Y)\rightarrow Y$. 
	
	In general:
	
	\begin{Def}
		If $f:Y\rightarrow X$ is any continuous map, then the \textbf{pullback} of $\eta:E\rightarrow X$ is 
		\begin{gather*}
			f^*\eta:f^*E\rightarrow Y
		\end{gather*}
		where $f^*E=Y\times_X E$
		\begin{center}
			\begin{tikzcd}						f^*E\arrow[r]\arrow[d,"f^*\eta"]&E\arrow[d,"\eta"]\\
				Y\arrow[r,"f"]&X
			\end{tikzcd}
		\end{center}
	\end{Def}
	
	\subsubsection{Vector Bundles}
	
	
	\begin{example}
		(Trivial Bundle) If $V$ is an $n$-dimensional vector space. The projection from $T^n=X\times V$ to $X$ is a \textbf{constant family}
		\begin{gather*}
			\eta:X\times V\rightarrow X
		\end{gather*}
		Such bundle is called \textbf{trivial vector bundle}.
	\end{example}
	
	\begin{Def}
		(Vector Bundle) A \textbf{vector bundle} is a family of vector spaces such that: 
		
		\begin{itemize}
			\item Local Trivial: For every point $x\in X$, $\eta|U$ is a trivial.
			\begin{center}
				\begin{tikzcd}
					\eta^{-1}(U)\arrow[rr,shift left=2]\arrow[rd]&&T=U\times V\arrow[ld]\arrow[ll,shift left=2]\\
					&U
				\end{tikzcd}
			\end{center}
		\end{itemize}
	\end{Def}
	
	So sometimes a vector bundle is called a \textbf{locally trivial} family of vectors.
	
	\begin{Def}
		(Category of vector bundles) Define the category $Vect_F(X)$, with the elements:
		
		\begin{itemize}
			\item $Ob(Vect_F(X))=\{$Vector Bundles on $X\}$
			\item $Mor(Vect_F(X))=\{$Morphisms between elements in $Ob(Vect_F(X))\}$
		\end{itemize}
		
		And a functor is given by: $f^*:Vect(Y)\rightarrow Vect(X)$.
	\end{Def}
	
	\subsubsection{Product}
	
	\begin{Def}
		Let $\eta:E\rightarrow X,\varphi:F\rightarrow X$ be two vector bundle. Then, the \textbf{product of vector bundles} is defined as $(E\times F,X,\eta\times \varphi)$
		
		\begin{gather*}
			\eta\times \varphi:E\times F\rightarrow X\times X
		\end{gather*}
	\end{Def}
	
	This gives a bilinear functor
	\begin{gather*}
		\boxplus:Vect(X)\times Vect(X)\rightarrow Vect(X\times X) 
	\end{gather*}
	
	
	
	\subsubsection{Tensor Product}
	
	\begin{Def}
		Let $\eta:E\rightarrow X,\varphi:F\rightarrow X$ be two vector bundles. Then,
		
		\begin{gather*}
			E\otimes F:=\coprod_{x\in X} E_x\otimes F_x
		\end{gather*}
		
		with $\eta\otimes \varphi:E\otimes F\rightarrow X,e\otimes f\mapsto x;e\in E_x,f\in F_x$.
	\end{Def}
	This is also a bilinear functor
	\begin{gather*}
		\otimes:Vect(X)\times Vect(X)\rightarrow Vect(X)
	\end{gather*}
	
	\subsubsection{Whitney Sum}
	
	\begin{Def}
		The \textbf{Whitney sum} $E\oplus F$ of two vector bundle $\eta:E\rightarrow X,\varphi:F\rightarrow X$ is the family of all vector spaces $E_x\oplus F_x$, which can be seen as a topological subspace of $E\times F$:
		\begin{gather*}
			(E_x\times F_x)_{x\in X}\hookrightarrow E\times F
		\end{gather*}
		With the new vector bundle construction induced by $\eta\times \varphi$:
		\begin{gather*}
			\eta\oplus \varphi:=\eta\times \varphi|E\oplus F
		\end{gather*}
	\end{Def}
	
	Whitney sum is the categoricial coporudct in $Vect(X)$. 
	\begin{center}
		\begin{tikzcd}
			&E\oplus F\arrow[d,dashrightarrow]\arrow[dd,bend right=50]\\
			E\arrow[rd]\arrow[ru]\arrow[r]&G\arrow[d]&F\arrow[ld]\arrow[lu]\arrow[l]\\
			&X
		\end{tikzcd}
	\end{center}
	
	Given $E\oplus F\rightarrow G:E_x\times F_x\ni (e,f)\mapsto i_1(e)+i_2(f)$ which is well-defined. 
	
	And commutes: Via the product structure. 
	
	Especially: $E\oplus F$ is also the categoricial coproduct:
	
	\begin{center}
		\begin{tikzcd}
			&G\arrow[ld]\arrow[d,dashrightarrow]\arrow[rd]\arrow[dd,bend right=50]\\
			E\arrow[rd]&E\oplus F\arrow[l]\arrow[r]\arrow[d]&F\arrow[ld]\\
			&X
		\end{tikzcd}
	\end{center}
	
	$G\rightarrow E\oplus F,g\mapsto p_1(g)+ p_2(g)$. Omitted.
	
	\begin{corollary}
		$Vect(X)$ is an additional category.
	\end{corollary}
	\begin{proof}
		Via checking definition.
	\end{proof}
	
	\begin{example}
		(Mobius Bundles) The Mobius bundle $E^1$ over $S$
		
		\begin{gather*}
			\pi:E^1\rightarrow S
		\end{gather*}
		With strip $E_x\cong \RR,x\in S^1$. 
		
		$E^1$ can be seen as the pasting of $[0,1]\times \RR$ via
		\begin{gather*}
			(0,v)\sim (1,-v)
		\end{gather*}
	\end{example}
	
	\begin{example}
		(Whitney Sum of Tangent Bundle of Product Manifolds) For the product of manifolds $M_1\times M_2$, we have
		
		\begin{gather*}
			T_{(x,y)}(M\times N)\cong T_x M\oplus T_y N
		\end{gather*}
		
		Thus, we have
		\begin{gather*}
			T(M\times N)=TM\oplus TN\rightarrow M\times N
		\end{gather*}
	\end{example}
	
	New Vector Bundles from Olds:
	
	\begin{itemize}
		\item Restsriction;
		\item Induced;
		\item Carstian Product
		\item Whitney Sum;
	\end{itemize}
	
	\subsubsection{Diemension of Vector Bundles}
	
	Notice that: $dim(E):X\rightarrow \NN,x\mapsto E_x$ is locally constant. 
	
	\begin{Def}
		An $n$-dimensional vector space is a bundle $E$ such that $dim(E)=n$ is a constant function. 
		
		Especially, on connected $X$, the function is constant.
	\end{Def}
	
	\begin{itemize}
		\item 
	\end{itemize}
	
	\begin{Def}
		内容...
	\end{Def}
	
	\begin{Def}
		Define:
		
		\begin{gather*}
			V_{n,F}(X):=
		\end{gather*}
	\end{Def}
	
	\subsubsection{Subbundles}
	
	\begin{Def}
		\begin{itemize}
			\item A \textbf{subbundle} of a vector space $\eta:E\rightarrow X$ is a subspace $F$ of $E$ which is a vector bundle under induced structure.
			
			i.e. $F_x$ is a vector subspace of $E_x$, and family $F\rightarrow X$ is locally trivial. 
			
			\item The \textbf{quotient bundle} $E/F$ is defined as the union of vector spaces
			\begin{gather*}
				(E_x/F_x)
			\end{gather*}
			given quotient topology(by $\eta$). 
		\end{itemize}
	\end{Def}
	
	Short exact sequence
	\begin{center}
		\begin{tikzcd}
			0\arrow[r]&F\arrow[r]&E\arrow[r]&E/F\arrow[r]&0
		\end{tikzcd}
	\end{center}
	
	in $VB(X)$.
	
	\begin{Def}
		A \textbf{vector bundle} $E\rightarrow X$ is said to be of \textbf{finite type} is there is a finite covering $U_1,\dots,U_n$ such that $E|U_i$ is a trivial bundle.
	\end{Def}
	
	\begin{proposition}
		\begin{itemize}
			\item 
			\item 
		\end{itemize}
	\end{proposition}
	
	\begin{proof}
		内容...
	\end{proof}
	
	\begin{theorem}
		(Subbundle Theorem) Let $E\rightarrow X$ be a vector bundle on \underline{paracomapct space} $X$:
		
		\begin{itemize}
			\item 
			\item 
		\end{itemize}
	\end{theorem}
	
	\begin{proof}
		With the conclusion \ref{Algebraic-K-Theory-Vector-Bundles-Metrics-on-Paracompact-Space}, there exists a metric on the bundle. 
	\end{proof}
	\subsubsection{Sections}
	
	\begin{Def}
		\begin{itemize}
			\item Let $\eta:E\rightarrow X$ be a vector space. A \textbf{global section} of $\eta$ is a continuous map $s:X\rightarrow E$ such that
			\begin{gather*}
				\eta\circ s=1_X
			\end{gather*}
			
			\item It is nowhere zero if $s(x)\not=0,\forall x\in X$.
		\end{itemize}
	\end{Def}
	In general, we could define the \textbf{local section} 
	\begin{gather*}
		s:U\rightarrow E\qquad\eta\circ s=id_U
	\end{gather*}
	
	\subsubsection{Vector Bundle under Riemannian Metrics}
	
	Review:
	
	\begin{itemize}
		\item Riemannian Metrics;
		\item Hermitian Metrics;
	\end{itemize}
	
	Metric on $E_x\Rightarrow$ Metric on $E$.
	
	A fundamental result is:
	
	\begin{theorem}
		\label{Algebraic-K-Theory-Vector-Bundles-Metrics-on-Paracompact-Space}If $X$ is a \underline{paracompact space}, then there exists a metric on $E$, when $F=\RR,\CC,\mathbb{H}$.
	\end{theorem}
	\begin{proof}
		内容...
	\end{proof}
	
	\subsubsection{Patching Vector Bundles}
	
	{\LARGE T}ransition of vector bundles: Consider two locally trivial maps
	\begin{gather*}
		内容...
	\end{gather*}
	
	\begin{Def}
		(Tensor Product of Bundles)
	\end{Def}
	
	\begin{Def}
		(Determinant Bundle)
	\end{Def}
	
	\subsubsection{Cancellation Theorems}
	
	\begin{theorem}
		(Real Cancellation Theorm)
	\end{theorem}
	
	\begin{proof}
		内容...
	\end{proof}
	
	\subsubsection{Homotopy Invariance Theory}
	
	\subsubsection{Classifying Vector Bundles}
	
	\begin{theorem}
		(Classifying Theorem) Let $X$ be paracompact space. Then, 
		
		\begin{gather*}
			VB_n(X)\cong [X,Grass_n]
		\end{gather*}
	\end{theorem}
	
	\begin{proof}
		内容...
	\end{proof}
	
	\begin{Def}
		(Classifying Spaces)
	\end{Def}
	
	\begin{theorem}
		(Classifying Theorem for Line Bundles) 
		
		\begin{gather*}
			内容...
		\end{gather*}
	\end{theorem}
	
	\begin{proof}
		内容...
	\end{proof}
	
	\subsection{Serre-Swan Theorem}
	
	\begin{proposition}
		(Global Section) Let $\eta:E\rightarrow X$ be a vector bundle. Let $\Gamma(E)$ be the set of global sections of $E$. 
		
		Then, $\Gamma(E)$ is a module over $C^0(X)$ of continuous functions on $X$. 
		
		Especially, if $E$ is a trivial bundle of dimension $n$, then $\Gamma(E)$ is a free $C^0(X)$-module of rank $n$.
	\end{proposition}
	
	\begin{proof}
		Module axiom:
		\begin{itemize}
			\item Sum:
			\item Scalar Product:
		\end{itemize}
	\end{proof}
	
	\begin{corollary}
		If $X$ is paracompact, then $\Gamma(E)$ is a locally free $C^0(X)$-module. 
		
		$\Gamma(E)$ is a finitely generated projective module if $X$ is compact or $E$ is of finite type.
	\end{corollary}
	
	\begin{proof}
		内容...
	\end{proof}
	
	\begin{theorem}
		(Swan) Let $X$ be a compact Hausdorff space, and write $R$ for $C^0(X)$. 
		
		Then, $\Gamma:VB(X)\rightarrow P(R)$ of finitely generated projective modules is a functor. The homomorphism
		\begin{gather*}
			\Gamma_{VB(X)} (E,F)
		\end{gather*}
		This implies, there exists an equivalence of categories $VB(X)\cong P(C^0(X))$. 
	\end{theorem}
	\begin{proof}
		内容...
	\end{proof}
	
	\textbf{Remark} The Serre's theorem will be given in AG.
	
	Here:
	\begin{itemize}
		\item Construction:
		\item Affines:
	\end{itemize}
	
	\subsubsection{Projective Flag Modules}
	
	\subsubsection{Splitting Principle}
	
	\subsubsection{Pontrijagin Classes}
	\subsection{Characteristic Classes}
	
	\textbf{Abstract} The determinant line bundle $det(E)$ of a vector bundle $E$ yields to a cohomology class. 
	
	\begin{itemize}
		\item 
		\item 
	\end{itemize}
	
	\subsubsection{Axioms for Stiefel-Whitney Classes}
	
	\begin{axiom}
		内容...
	\end{axiom}
	
	\subsubsection{Axioms for Chern Classes}
	
	\begin{axiom}
		内容...
	\end{axiom}
	\subsubsection{Pontrjagin Classes}
	
	\begin{Def}
		(Complex Conjugate Bundle) 
	\end{Def}
	
	\subsection{Others}
	\subsubsection{Complexification of Real Vector Bundlees}
	\section{Algebraic Vector Bundles}
\end{document}